\documentclass[10pt,a4paper]{amsart}
\usepackage[margin=19mm]{geometry}
\usepackage{amsmath,amssymb,mathtools}
\usepackage[scr=rsfs,cal=euler]{mathalfa}
\usepackage{xfrac}
\usepackage[unicode,hidelinks,bookmarks=false]{hyperref}
\newcommand{\ie}{i.e.\ }
\newcommand{\cf}{cf.\ }
\newcommand{\resp}{respectively\ }
\newcommand{\Subsection}{\subsection}
\providecommand{\nobreakdash}{\nobreak\hbox{-}}
\setlength{\emergencystretch}{2em}
\multlinegap0pt
\usepackage{amssymb}
\usepackage{xcolor}
\newtheorem{lemma}{Lemma}
\newtheorem{prop}{Proposition}
\usepackage{cleveref}
\crefname{lemma}{Lemma}{Lemmas}
\crefname{prop}{Proposition}{Propositions}
\newcommand{\B}[2]{\langle#1,#2\rangle}
\newcommand{\CS}{\operatorname{CS}}
\newcommand\Cal{\mathcal}
\newcommand{\Om}{\Omega}
\newcommand*{\R}{\mathbb{R}}
\newcommand{\Z}{\mathbb{Z}}
\newcommand{\col}{:}
\renewcommand*{\d}{\mathrm{d}}
\newcommand{\deq}{:=}
\renewcommand\frak{\mathfrak}
\renewcommand{\th}{\theta}
\newcommand{\tr}{\operatorname{tr}}
\newcommand{\x}{\times}
\title{Flat extensions of principal connections and the Chern--Simons $\boldsymbol{3}$-form}
\author{Andreas {\v{C}}ap, Keegan J.~Flood, Thomas Mettler}
\date{Source: arXiv:2409.12811v4 (CC BY 4.0)}
\begin{document}
\maketitle

\noindent\emph{Source and licence.} Flat extensions of principal connections and the Chern--Simons $\boldsymbol{3}$-form. Version arXiv:2409.12811v4 (CC BY 4.0), distributed by arXiv under the Creative Commons Attribution 4.0 International licence, http://creativecommons.org/licenses/by/4.0/, as recorded in the arXiv licence metadata for this exact version and shown on the abstract page https://arxiv.org/abs/2409.12811v4. Full licence notice: \texttt{licenses/arxiv-2409.12811-CC-BY-4.0.txt}.\par\smallskip

\noindent\textit{Editorial scope.} Counted unit: the article's prop prop:normalization (source0.tex lines 919--928). The article's complete original proof of it is delivered with it (source0.tex lines 929--1001, one \texttt{proof} environment of 73 lines). No mathematics is written, repaired or re-derived: the statement and the proof are the article's own lines, in the article's own notation, and no retained line is substituted or re-flowed.  Retained uncounted as the source-local context the proof needs: lines 436--445; lines 486--494.  Nothing was omitted from any retained span, so the package carries no omission record.  No retained line contains a citation, so the wrapper carries no bibliography and no bibliography entry.  No retained line contains a figure environment, an image-inclusion command or drawing code, so no image is delivered and the package declares no asset dependency.  Editorial formatting only, in addition: an AMS \texttt{amsart} preamble that restates the article's own declarations and macros used by the retained lines, namely the environments \texttt{lemma}, \texttt{prop} on separate counters, together with the standard packages those lines need.  The retained environments print in this document as Lemma 1, Proposition 1 in order of appearance, whereas the article prints the same items with the numbering of its own full document.  The complete native source of the article as distributed in the arXiv e-print for this exact version is bundled unchanged as \texttt{sources/165511/source0.tex}.
\begin{lemma}\label{lem:CS-pairs}
Let $\th\in\Om^1(N,\mathfrak{\tilde{g}})$ be such that for the decomposition
$\th=\th^\top+\th^\perp$, we have
$[\th^\perp,\th^\perp]\in\Om^2(N,\mathfrak{g})$. Then for
$\Theta=\d\th+\frac12[\th,\th]$ with associated decomposition
$\Theta=\Theta^\top+\Theta^\perp$, we obtain
\begin{equation}\label{eqn:partialblindness}
\CS(\th)=\CS(\th^\top)+\B{\th^\perp}{\Theta^\perp}. 
\end{equation}
\end{lemma}

\begin{lemma}\label{lem:CS-subgroup}
Let $i\col G\to\tilde G$ be an inclusion of a Lie subgroup and consider the corresponding
subalgebra $\frak{g}\subset\tilde{\frak{g}}$. Assume that $\B{\,\cdot\,}{\cdot\,}$ is an
invariant bilinear form on $\tilde{\frak{g}}$ such that $\B{\,\cdot\,}{\cdot\,}|_{\frak{g}\x\frak
g}$ is non-degenerate.

If $\B{\,\cdot\,}{\cdot\,}$ is normalized in such a way that $[\CS(\mu_{\tilde G})]\in
H^3(\tilde G,\Z)$, then $[\CS^{\frak{g}}(\mu_G)]\in H^3(G,\Z)$.
\end{lemma}

\begin{prop}\label{prop:normalization}
Let $\mu_4$ denote the Maurer--Cartan form of $\mathrm{SO}(4)$. Then
\begin{equation}\label{eqn:goodnorm}
\zeta\deq\frac{1}{16\pi^2}\tr\left(\mu_4\wedge \d
\mu_4+\tfrac{2}{3}\mu_4\wedge\mu_4\wedge\mu_4\right)
\end{equation}
represents an element of $H^3(\mathrm{SO}(4),\Z)$.  Moreover, for the
  inclusion $\iota:\mathrm{SO}(3)\to \mathrm{SO}(4)$, we have
  $\int_{\mathrm{SO}(3)}\iota^*\zeta=1$.
\end{prop}

\begin{proof}
We will show that the $3$-form defined in \eqref{eqn:goodnorm} evaluated on a set of
generators of $H_3(\mathrm{SO}(4),\Z)$ is $1$.  The principal right
$\mathrm{SO}(3)$-bundle $\pi \col \mathrm{SO}(4) \to S^3$ is trivial so that
$\mathrm{SO}(4)$ is diffeomorphic to $S^3\times \mathrm{SO}(3)$. Standard results
about the homology of $S^3$ and $\mathrm{SO}(3)\simeq \mathbb{RP}^3$ together with
K\"unneth's theorem imply that
\[
H_3(\mathrm{SO}(4),\Z)\simeq H_3(S^3\times \mathrm{SO}(3),\Z)\simeq \Z \oplus \Z. 
\]
Moreover, $H_3(\mathrm{SO}(4),\Z)$ is generated by the pushforward of the fundamental
class of $\mathrm{SO}(3)$ under the inclusion $\iota \col \mathrm{SO}(3) \to
\mathrm{SO}(4)$ and the pushforward of the fundamental class of $S^3$ under a section
$\sigma \col S^3 \to \mathrm{SO}(4)$. It is thus sufficient to show that we have
\begin{equation}\label{eqn:integrals}
\int_{\mathrm{SO}(3)}\iota^*\zeta=1 \qquad \text{and}\qquad \int_{S^3}\sigma^*\zeta=1
\end{equation}
with respect to suitable orientations of $S^3$ and $\mathrm{SO}(3)$. We first treat
the left integral.  We know from the proof of \cref{lem:CS-subgroup}
  that $\iota^*\mu_4^\top=\mu_3$, the Maurer--Cartan form of
$\mathrm{SO}(3)$. Furthermore, $(\tilde{\mathfrak{g}},\mathfrak{g})$ is a symmetric
pair so that \cref{lem:CS-pairs} implies
\[
\int_{\mathrm{SO}(3)}\iota^*\zeta=\frac{1}{16\pi^2}\int_{\mathrm{SO}(3)}\tr\left(\mu_3\wedge\d\mu_3+\tfrac{2}{3}\mu_3\wedge\mu_3\wedge\mu_3\right).
\]
Writing
\[
\mu_3=\begin{pmatrix} 0 & -\omega_1 & -\omega_2 \\ \omega_1 & 0 & -\psi \\ \omega_2 & \psi & 0 \end{pmatrix}
\]
for left-invariant $1$-forms $\omega_1,\omega_2,\psi \in \Omega^1(\mathrm{SO}(3))$, an elementary computation gives
\[
\tr\left(\mu_3\wedge\d\mu_3+\tfrac{2}{3}\mu_3\wedge\mu_3\wedge\mu_3\right)=2\,\omega_1\wedge\omega_2\wedge\psi.
\]
Under the identification $\mathrm{SO}(3)\cong\Cal{SO}S^2$ observed above, the integral 
\[
\int_{\mathrm{SO}(3)}\omega_1\wedge\omega_2\wedge\psi
\] 
equals the product of the length of the typical fibre of $\mathrm{SO}(2)\to \mathcal{SO}S^2 \to S^2$ with the surface area of $S^2$. The former is given by $2\pi$ and the latter by $4\pi$. In summary we thus have
\[
\int_{\mathrm{SO}(3)}\iota^*\zeta=\frac{2}{16\pi^2}\int_{\mathrm{SO}(3)}\omega_1\wedge\omega_2\wedge\psi=\frac{1}{8\pi^2}\int_{\mathrm{SO}(3)}\omega_1\wedge\omega_2\wedge\psi=1. 
\]

In order to compute the right integral of \eqref{eqn:integrals} we consider the smooth $\pi$-section $\sigma \col S^3 \to \mathrm{SO}(4)$ given by the rule
\[
x=\begin{pmatrix} x_1 \\ x_2 \\ x_3 \\ x_4 \end{pmatrix} \mapsto \begin{pmatrix} x_1 & -x_2 & -x_3 & -x_4 \\ x_2 & x_1 & x_4 & -x_3 \\ x_3 & -x_4 & x_1 & x_2 \\ x_4 & x_3 & -x_2 & x_1 \end{pmatrix}
\]
for all $x \in S^3$. A tedious but straightforward calculation shows that
\[
\sigma^*\mu_4=\begin{pmatrix} 0 & \kappa & -\xi & -\rho  \\ -\kappa & 0 & \rho & -\xi \\ \xi & -\rho & 0 & -\kappa \\ \rho & \xi & \kappa & 0 \end{pmatrix}
\]
where $\xi,\rho,\kappa \in \Omega^1(S^3)$ are given by
\begin{equation}
\begin{aligned}\label{eqn:leftinvariantforms}
\xi&=-x_3\d x_1+x_4\d x_2+x_1\d x_3-x_2 \d x_4,\\
\rho&=-x_4\d x_1-x_3\d x_2+x_2\d x_3+x_1\d x_4,\\
\kappa&=x_2\d x_1-x_1\d x_2+x_4\d x_3-x_3 \d x_4.
\end{aligned}
\end{equation}
From this one computes
\[
\sigma^*\zeta=\frac{8}{16\pi^2}\,\xi\wedge\rho\wedge\kappa=\frac{1}{2\pi^2}\,\xi\wedge\rho\wedge\kappa
\]
and
\begin{multline}\label{eqn:volformthreesphere}
\xi\wedge\rho\wedge \kappa=x_4\d x_1\wedge \d x_2 \wedge \d x_3-x_3\d x_1\wedge\d x_2\wedge \d x_4\\
+x_2\d x_1\wedge \d x_3\wedge \d x_4-x_1\d x_2\wedge \d x_3 \wedge\d x_4,
\end{multline}
where we use that $|x|^2=1$. The right hand side of \eqref{eqn:volformthreesphere} equals the restriction to $S^3$ of the interior product of the inwards pointing radial vector field $-\sum_{i}x_i\partial_i$ and the standard volume form $\d x^1\wedge\d x^2\wedge \d x^3\wedge \d x^4$ of $\R^4$. This is the standard volume form of $S^3$ with respect to the inwards orientation. We thus have
\begin{equation}\label{eqn:volumethresphere}
\int_{S^3}\xi\wedge \rho\wedge \kappa=2\pi^2
\end{equation}
and the claim follows.
\end{proof}

\end{document}
